\documentclass[11pt,a4paper]{article}

\usepackage[a4paper,margin=1in]{geometry}
\usepackage{fontspec}
\IfFontExistsTF{Times New Roman}{\setmainfont{Times New Roman}}{\setmainfont{TeX Gyre Termes}}
\IfFontExistsTF{Menlo}{\setmonofont{Menlo}}{\setmonofont{Latin Modern Mono}}
\usepackage{amsmath,amssymb}
\usepackage{booktabs}
\usepackage{longtable}
\usepackage{array}
\usepackage{enumitem}
\usepackage{graphicx}
\usepackage{fancyhdr}
\usepackage[round,authoryear]{natbib}
\usepackage[hidelinks]{hyperref}
\usepackage{xurl}
\usepackage{microtype}

\setlength{\parindent}{0pt}
\setlength{\parskip}{0.65em}
\setlist[itemize]{leftmargin=1.4em}
\newcommand{\file}[1]{\path{#1}}
\newcommand{\code}[1]{\path{#1}}
\newcolumntype{L}[1]{>{\raggedright\arraybackslash}p{#1}}

\pagestyle{fancy}
\fancyhf{}
\lhead{\small Quantitative Risk Analytics}
\rhead{\small Controlled Expansion After Pilot 01}
\cfoot{\thepage}
\renewcommand{\headrulewidth}{0.4pt}

\title{\textbf{Quantitative Risk Analytics Report}\\
\large Small Controlled Expansion After Pilot 01 Stress Maps}
\author{Codex-assisted diagnostic analytics report}
\date{June 17, 2026}

\begin{document}
\maketitle

\begin{abstract}
This report documents a small quantitative risk-analytics expansion after Pilot 01. The experiment
uses the validated baseline American Crank--Nicolson / projected SOR solver only. It quantifies
boundary movement, continuation premium, Gamma concentration, PSOR effort, runtime, and LCP
diagnostics under selected one-at-a-time stress parameters. The outputs are diagnostic analytics
only. They are not broad stress-map datasets, training labels, neural-surrogate inputs, or
production risk surfaces.
\end{abstract}

\tableofcontents
\newpage

\section{Purpose of Quantitative Risk Analytics}

Pilot 01 showed that the validated baseline solver can produce small, metadata-rich diagnostic
stress-map figures. The present step is narrower than data generation but more quantitative than
visual inspection. Its purpose is to summarize how key numerical and financial diagnostics move
when selected stress parameters are changed one at a time.

The experiment follows the downstream design gate and Ticket 12 synthesis decision,
\code{PASS_SOLVER_VALIDATION_WITH_LIMITATIONS}. It keeps the same classical finite-difference
American option setting described by the project validation ladder
\citep{black_scholes_1973,merton_1973,brennan_schwartz_1977,wilmott_howison_dewynne_1995}.

\section{Link to Pilot 01 and Senior Handout Experiment 04}

The senior handout's risk-analytics stage asks for boundary movement, Gamma concentration, PSOR
iteration effort, runtime, and complementarity diagnostics before dataset construction. This
report implements that idea at small scale. It deliberately does not export training arrays or
machine-learning-ready targets.

\begin{longtable}{L{0.30\textwidth} L{0.54\textwidth}}
\toprule
Handout theme & This report's controlled implementation \\
\midrule
\endhead
Boundary movement & Selected boundary values and boundary curves across stress sweeps. \\
Gamma concentration & Full-grid and strict-mask maximum \(|\Gamma|\) summaries. \\
PSOR effort & Mean and maximum PSOR iterations plus runtime seconds. \\
Complementarity diagnostics & Obstacle, equation, and complementarity product metrics. \\
Dataset caution & Outputs are marked \code{analytics\_diagnostic\_only}. \\
\bottomrule
\end{longtable}

\section{Solver and Grid Choice}

All runs use the Ticket 04 baseline American CN/PSOR solver:
\begin{center}
\code{american\_crank\_nicolson\_psor\_price}
\end{center}
Rannacher smoothing is not used. Ticket 10A kept Rannacher as an optional comparison tool only,
so this analytics step does not silently change the validated solver.

Every run uses normalized strike \(K=1\), \(T=1\), \(S_{\max}=4\), and \(M=N=120\). The main
interpretation region remains \(S/K \in [0.4,1.8]\), while the solver still computes on the full
domain.

\section{Stress Parameter Design}

\begin{longtable}{L{0.23\textwidth} L{0.25\textwidth} L{0.37\textwidth}}
\toprule
Family & Sweep & Values \\
\midrule
\endhead
American put & Volatility & \(\sigma \in \{0.20,0.40,0.60\}\), with \(r=0.05\), \(q=0.02\). \\
American put & Dividend yield & \(q \in \{0.00,0.02,0.04\}\), with \(r=0.05\), \(\sigma=0.20\). \\
American put & Interest rate & \(r \in \{0.03,0.05,0.08\}\), with \(q=0.02\), \(\sigma=0.20\). \\
Dividend call & Dividend yield & \(q \in \{0.03,0.08,0.10,0.14\}\), with \(r=0.05\), \(\sigma=0.20\). \\
Dividend call & Volatility & \(\sigma \in \{0.20,0.40,0.60\}\), with \(r=0.05\), \(q=0.08\). \\
Dividend call & Interest rate & \(r \in \{0.03,0.05,0.08\}\), with \(q=0.08\), \(\sigma=0.20\). \\
Dividend call & Small \(q\)-\(\sigma\) diagnostic & \(q \in \{0.03,0.08,0.10,0.14\}\) and \(\sigma \in \{0.20,0.40,0.60\}\). \\
\bottomrule
\end{longtable}

The experiment writes 31 analytics rows: 19 one-at-a-time sweep rows and 12 small \(q\)-\(\sigma\)
diagnostic heatmap rows. Duplicate parameter solves are cached internally, but the CSV rows keep
their sweep membership.

\section{Metric Definitions}

\begin{longtable}{L{0.26\textwidth} L{0.56\textwidth}}
\toprule
Metric & Meaning \\
\midrule
\endhead
Boundary moneyness & Threshold-based boundary spot divided by \(K\), extracted from \(U-\Phi\). \\
Continuation premium & \(U-\Phi\), the value above immediate exercise payoff. \\
PSOR effort & Maximum and mean projected-SOR iteration counts over time steps. \\
Runtime seconds & Wall-clock runtime for the baseline solver and attached diagnostics. \\
Obstacle violation & \(\max(\Phi-U,0)\). \\
Equation violation & \(\max(b-AU,0)\) from the reconstructed LCP. \\
Complementarity product & \((U-\Phi)(AU-b)\), summarized by maximum absolute value. \\
Full max \(|\Gamma|\) & Maximum absolute finite-difference Gamma over finite grid entries. \\
Strict max \(|\Gamma|\) & Maximum absolute Gamma after excluding boundary-near, payoff-kink, and maturity rows. \\
\bottomrule
\end{longtable}

\section{Tables and Interpretation}

\subsection{Boundary Direction Checks}

The main qualitative handout directions are visible in the generated numbers. At \(\tau/T=1\), the
American put boundary falls as volatility rises:
\[
0.7668 \rightarrow 0.5667 \rightarrow 0.4000
\]
for \(\sigma=0.20,0.40,0.60\). This matches the expected direction: higher volatility raises the
value of waiting, so the put exercise boundary moves lower.

For the dividend-paying call, the boundary also falls as dividend yield rises:
\[
1.8979 \rightarrow 1.2666 \rightarrow 1.2332 \rightarrow 1.1666
\]
for \(q=0.03,0.08,0.10,0.14\). This matches the expected direction that higher dividend yield
makes early exercise more attractive for calls.

\begin{longtable}{L{0.23\textwidth} L{0.30\textwidth} L{0.31\textwidth}}
\toprule
Sweep & Boundary at \(\tau/T=1\) & Interpretation \\
\midrule
\endhead
Put \(\sigma=0.20,0.40,0.60\) & 0.7668, 0.5667, 0.4000. & Decreasing; matches expected volatility direction. \\
Call \(q=0.03,0.08,0.10,0.14\) & 1.8979, 1.2666, 1.2332, 1.1666. & Decreasing; matches expected dividend-yield direction. \\
Call \(\sigma=0.20,0.40,0.60\) & 1.2666, 1.7331, 2.3660. & Increasing; higher volatility delays high-spot exercise in this diagnostic. \\
\bottomrule
\end{longtable}

\subsection{PSOR, Runtime, and LCP Checks}

All 31 analytics rows converged and were marked acceptable for analytics review. The largest
equation violation was \(4.8938\times 10^{-8}\), and the largest complementarity product was
\(4.7325\times 10^{-10}\). These are small diagnostic values, but they are not a mathematical
proof across all future parameter settings.

Volatility increased PSOR effort and runtime in both families. For the put volatility sweep, mean
PSOR iterations rose from \(6.99\) to \(29.42\) as \(\sigma\) moved from \(0.20\) to \(0.60\). The
runtime for the same sweep rose from about \(0.049\) seconds to about \(0.204\) seconds. For the
call volatility sweep, mean iterations rose from \(7.35\) to \(14.92\).

\begin{longtable}{L{0.28\textwidth} L{0.20\textwidth} L{0.20\textwidth} L{0.20\textwidth}}
\toprule
Diagnostic & Minimum / baseline reading & Maximum observed & Comment \\
\midrule
\endhead
Obstacle violation & \(0\) & \(0\) & Payoff obstacle respected in all runs. \\
Equation violation & Near \(10^{-9}\) for base cases & \(4.8938\times10^{-8}\) & Largest value occurs in high-volatility put run. \\
Complementarity product & Around \(10^{-11}\) for many runs & \(4.7325\times10^{-10}\) & Still small relative to the diagnostic threshold. \\
Mean PSOR iterations & Around \(7\) in base cases & \(29.42\) & High volatility increases iteration effort. \\
Runtime seconds & Around \(0.05\) in base cases & \(0.204\) & Small absolute runtime in this pilot. \\
\bottomrule
\end{longtable}

\subsection{Gamma Concentration}

Full-grid maximum \(|\Gamma|\) is \(30\) across the runs because the full metric includes fragile
regions such as payoff kink and maturity-row effects. The strict-mask maximum is much smaller,
around \(1.92\) to \(3.29\) in the one-at-a-time sweeps. This difference is exactly why the project
keeps Gamma diagnostic and masked. The strict-mask Gamma values do not move monotonically in every
sweep, so the report does not claim a simple universal Gamma direction.

\section{Figures and Interpretation}

Figure~\ref{fig:put_boundary_scalar} shows the scalar put boundary direction check. It is the
cleanest evidence that the selected put boundary decreases with volatility in this pilot.

\begin{figure}[htbp]
\centering
\includegraphics[width=0.78\textwidth]{results/03_quantitative_risk_analytics/figures/risk_analytics_put_boundary_vs_volatility.png}
\caption{American put boundary moneyness versus volatility at \(\tau/T=1\).}
\label{fig:put_boundary_scalar}
\end{figure}

Figure~\ref{fig:call_boundary_scalar} shows the corresponding dividend-call direction check. The
boundary decreases as dividend yield rises.

\begin{figure}[htbp]
\centering
\includegraphics[width=0.78\textwidth]{results/03_quantitative_risk_analytics/figures/risk_analytics_call_boundary_vs_dividend_yield.png}
\caption{Dividend-call boundary moneyness versus dividend yield at \(\tau/T=1\).}
\label{fig:call_boundary_scalar}
\end{figure}

Figures~\ref{fig:put_boundary_curves} and~\ref{fig:call_boundary_curves} show the whole
threshold-extracted boundary curves over time-to-maturity. They are approximate diagnostic curves,
not exact analytical free-boundary solutions.

\begin{figure}[htbp]
\centering
\includegraphics[width=0.78\textwidth]{results/03_quantitative_risk_analytics/figures/risk_analytics_put_vol_boundary_curves.png}
\caption{American put boundary curves across the volatility sweep.}
\label{fig:put_boundary_curves}
\end{figure}

\begin{figure}[htbp]
\centering
\includegraphics[width=0.78\textwidth]{results/03_quantitative_risk_analytics/figures/risk_analytics_call_q_boundary_curves.png}
\caption{Dividend-call boundary curves across the dividend-yield sweep.}
\label{fig:call_boundary_curves}
\end{figure}

Figure~\ref{fig:gamma_concentration} separates full-grid Gamma from strict-mask Gamma. The full
values are dominated by fragile locations; the strict values are the safer summary for human
review, but still not production Greek labels.

\begin{figure}[htbp]
\centering
\includegraphics[width=0.90\textwidth]{results/03_quantitative_risk_analytics/figures/risk_analytics_gamma_concentration.png}
\caption{Gamma concentration across key stress sweeps.}
\label{fig:gamma_concentration}
\end{figure}

Figure~\ref{fig:psor_runtime} shows numerical workload. The high-volatility cases are visibly
more expensive in iteration count and runtime, matching the senior handout's warning that stress
regimes can increase solver effort.

\begin{figure}[htbp]
\centering
\includegraphics[width=0.95\textwidth]{results/03_quantitative_risk_analytics/figures/risk_analytics_psor_runtime_iterations.png}
\caption{PSOR iteration and runtime diagnostics across analytics cases.}
\label{fig:psor_runtime}
\end{figure}

Figure~\ref{fig:lcp_stability} shows LCP diagnostic magnitudes. The values remain small, with the
largest equation violations appearing in higher-volatility runs.

\begin{figure}[htbp]
\centering
\includegraphics[width=0.95\textwidth]{results/03_quantitative_risk_analytics/figures/risk_analytics_lcp_stability.png}
\caption{LCP diagnostic stability across analytics cases.}
\label{fig:lcp_stability}
\end{figure}

Figure~\ref{fig:call_heatmap} is the optional small dividend-call \(q\)-\(\sigma\) boundary
heatmap. It is an aggregate diagnostic visualization only. It is not a saved data array and not a
dataset construction step.

\begin{figure}[htbp]
\centering
\includegraphics[width=0.70\textwidth]{results/03_quantitative_risk_analytics/figures/risk_analytics_call_q_sigma_boundary_heatmap.png}
\caption{Small aggregate dividend-call boundary moneyness heatmap over \(q\) and \(\sigma\).}
\label{fig:call_heatmap}
\end{figure}

\section{Qualitative Comparison with Senior Handout}

\begin{longtable}{L{0.34\textwidth} L{0.18\textwidth} L{0.34\textwidth}}
\toprule
Expected direction or caution & Result & Comment \\
\midrule
\endhead
Put boundary should decrease as volatility rises. & Matches. & Boundary at \(\tau/T=1\) falls from \(0.7668\) to \(0.4000\). \\
Dividend-call boundary should decrease as dividend yield rises. & Matches. & Boundary falls from \(1.8979\) at \(q=0.03\) to \(1.1666\) at \(q=0.14\). \\
Higher volatility may increase PSOR effort and runtime. & Matches. & Put mean iterations rise to \(29.42\); runtime rises to \(0.204\) seconds. \\
Gamma concentration should be strongest near kink/boundary regions. & Caution confirmed. & Full max \(|\Gamma|\) is dominated by fragile full-grid locations; strict-mask values are lower. \\
\bottomrule
\end{longtable}

\section{Acceptance Checks}

\begin{longtable}{L{0.42\textwidth} L{0.18\textwidth} L{0.28\textwidth}}
\toprule
Criterion & Result & Interpretation \\
\midrule
\endhead
Use baseline CN/PSOR solver only & PASS & No Rannacher run is used. \\
Keep sweeps small and controlled & PASS & 19 one-at-a-time rows plus 12 small heatmap rows. \\
All PSOR steps converge or are reported & PASS & All rows converged. \\
Obstacle violation near zero & PASS & Maximum obstacle violation is \(0\). \\
LCP diagnostics near tolerance & PASS & Diagnostics remain small in all rows. \\
Boundary metadata present & PASS & Selected-tau boundary rows include found/not-found status. \\
Greek diagnostics masked and cautious & PASS & Greek table includes strict-mask and caution metrics. \\
No dataset or neural outputs & PASS & Only CSV summaries, PNG figures, and this report are created. \\
\bottomrule
\end{longtable}

\section{Limitations}

\begin{itemize}
  \item The \(q\)-\(\sigma\) heatmap is a small aggregate diagnostic, not a broad two-dimensional
  parameter sweep.
  \item Boundary extraction is threshold-based with threshold \(10^{-6}\).
  \item The scalar direction checks use the nearest \(\tau/T=1\) row and should not be read as
  exact comparative statics.
  \item Gamma remains finite-difference diagnostic output. It is not a production Greek measure.
  \item Runtime seconds are local-machine observations, useful for this pilot but not a portable
  performance benchmark.
  \item The experiment supports human review, not automated downstream generation.
\end{itemize}

\section{What This Analytics Step Supports}

This step supports the claim that the validated baseline solver can produce small quantitative
diagnostic analytics aligned with the senior handout's risk-analytics expectations. It also
supports the specific qualitative observations that the put boundary decreases with volatility,
the dividend-call boundary decreases with dividend yield, and high volatility increases PSOR
effort in the tested settings.

\section{What This Analytics Step Does Not Support}

This step does not support broad stress-map generation, exact free-boundary claims, production
Greek reporting, training-label export, neural-surrogate readiness, or final paper-style claims
about all parameter regimes. Those remain blocked until a separate human-reviewed plan approves
the next stage.

\section{Recommended Next Step}

The recommended next step is human review of the analytics tables, figures, and limitations. If
accepted, the next technical step should be another small approved pilot that either enlarges the
stress grid modestly or compares selected analytics at a higher grid resolution. It should still
avoid dataset construction and surrogate-model training.

\section{Reproducibility Record}

\begin{longtable}{L{0.34\textwidth} L{0.50\textwidth}}
\toprule
Item & Record \\
\midrule
\endhead
Source module & \file{src/american_risk_surfaces/downstream/risk_analytics.py}. \\
Experiment script & \file{experiments/12_quantitative_risk_analytics.py}. \\
Test file & \file{tests/test_quantitative_risk_analytics.py}. \\
Output folder & \file{results/03_quantitative_risk_analytics/}. \\
Report source & \file{reports/04_downstream/quantitative_risk_analytics_report.tex}. \\
Report PDF & \file{reports/04_downstream/quantitative_risk_analytics_report.pdf}. \\
Bibliography & \file{reports/03_solver/references.bib}. \\
\bottomrule
\end{longtable}

\bibliographystyle{plainnat}
\bibliography{reports/03_solver/references}

\end{document}
